% ============================================================================
% INTRODUCTION (Condensed)
% ============================================================================

\section{Introduction}
\label{sec:introduction}

The Large Language Model (LLM) landscape has shifted from a monolith dominated by a few frontier generalists to a fragmented ecosystem of specialized capabilities. Developers now face a ``long tail'' of options: from reasoning powerhouses (e.g., GPT-4o, Claude 3.5) to domain-tuned specialists (e.g., DeepSeek Coder) and ultra-low-cost distillations (e.g., Gemini Flash Lite). While this fragmentation offers the theoretical possibility of drastically reducing costs, operationalizing it remains an open challenge. The central question is no longer availability, but selection: \emph{for a given prompt, what is the cheapest model that satisfies the user's constraints?}

Prior approaches to this \emph{model routing} problem generally fall into two categories: static regressors (e.g., RouteLLM)~\cite{ong2024routellm} or cascading chains (e.g., FrugalGPT)~\cite{chen2023frugalgpt}. While effective for small pools ($N<5$), these systems fail to scale to the modern reality of massive registries ($N>80$). Cascading architectures suffer from linear latency scaling---checking models sequentially introduces unacceptable delays---and thus force practitioners to manually curate rigid, short chains, leaving most long-tail specialists unutilized.

In this work, we introduce \ours{}, a framework that transitions from reactive chaining to scalable predictive routing. Unlike cascades that verify quality ex-post (after paying for a cheap generation), \ours{} uses a contextual bandit~\cite{slivkins2019contextual} with \emph{shippable priors} to predict model suitability ex-ante. By learning prompt--difficulty structure offline and shipping it as compact covariance priors, the router can warm-start in new domains and estimate which model will succeed before any LLM call is made. This yields a constant number of model invocations per request (unless verification triggers) while retaining access to large model pools.

We further address operational rigidity by unifying cost optimization and reliability into a tiered architecture. While purely predictive routing establishes a low-cost operating regime, we identify a ``confident failure'' mode in complex instruction-following tasks where reactive verifiers can be fooled by plausible-but-wrong outputs. To mitigate this, we propose a hybrid architecture that uses prediction as a gatekeeper for verification, dynamically trading cost for assurance via a tunable SLA parameter.

\paragraph{Contributions.}
\begin{itemize}[leftmargin=*]
    \item \textbf{Shippable Priors for Cold-Start Routing:} we distill offline teacher supervision into compact covariance priors. This enables few-shot domain calibration and reduces day-one regret relative to naive bandits (\S~\ref{sec:distillation}, \S~\ref{sec:rq1}).
    \item \textbf{Zero-Overhead Scalability:} we decouple prompt-difficulty learning from model-specific benchmarking. New models can be onboarded in constant time using public metadata as coarse safety filters; the bandit then autonomously refines its estimates via online exploration, removing the offline profiling bottleneck (\S~\ref{sec:scalability}).
    \item \textbf{Tiered \& Tunable Architecture:} we unify predictive routing and selective verification into a single framework. Hybrid mode improves instruction-following reliability by mitigating verifier blindness, outperforming reactive verification on this domain (98\% vs 95\%) (\S~\ref{sec:hybrid}, \S~\ref{sec:confident_failure}).
    \item \textbf{Empirical Benchmarking on Real Data:} evaluating on a registry of 81 models using real-world pricing and performance traces, we show that \ours{} defines a new Pareto frontier and enables smooth control over the cost--accuracy trade-off via a calibrated cascade threshold (\S~\ref{sec:evaluation}).
\end{itemize}
